\textbf{结论名称：} 对立事件概率公式（正难则反）

\textbf{适用条件：} 任意随机试验中，事件 \(A\) 及其对立事件 \(\overline{A}\) 均已定义，样本空间明确。

\textbf{变量说明：} \(A\) 为所关心的事件，\(\overline{A}\) 为 \(A\) 的对立事件（即 \(A\) 不发生），\(P(A)\) 与 \(P(\overline{A})\) 分别表示二者的概率，且满足 \(P(A)+P(\overline{A})=1\)。

\textbf{核心公式：} 
\[
\boxed{P(\overline{A}) = 1 - P(A)}
\]
等价地 \(P(A) = 1 - P(\overline{A})\)。

\textbf{使用场景：} 当直接计算事件 \(A\) 的概率步骤较多或情况复杂，而其对立事件 \(\overline{A}\) 的概率更易求得时，运用“正难则反”的策略。常见于“至少有一个”“至多”“不小于”等表述的概率问题。

\textbf{简单例子：} 掷一枚均匀骰子，求点数大于 \(2\) 的概率。设事件 \(A = \{\text{点数}\le 2\} = \{1,2\}\)，则 \(P(A)=\frac{2}{6}=\frac{1}{3}\)。所求为对立事件 \(\overline{A}=\{\text{点数}>2\}\) 的概率，由公式得 \(P(\overline{A}) = 1 - \frac{1}{3} = \frac{2}{3}\)，直接验证亦得 \(4/6=2/3\)。